% =============================================================================
% Chapter 6 — Conclusion
% =============================================================================

\chapter{Conclusion}
\label{chap:conclusion}

\index{conclusion}

\section{Summary of Contributions}

\index{contributions!summary}

This work presented \textbf{QRC-EV}, a complete framework for quantum
reservoir computing applied to electric vehicle charging load forecasting.
We demonstrated that the combination of a GPU-accelerated quantum reservoir
with a stateful Quantum Hidden Markov Model enables competitive performance
on temporal learning tasks, particularly those requiring explicit memory
of long histories.

\index{contributions!QRC-EV framework}
\index{contributions!QHMM-OOM}
\index{contributions!TD(λ)}

The specific technical contributions are:

\index{contributions}
\begin{enumerate}
  \item \textbf{Stateful GPU quantum reservoir.} We implemented a
    CUDA Quantum-based reservoir with input scaling, sparse ZZ gates,
    and depolarizing channel state maintenance, scalable to 30 qubits
    on NVIDIA Blackwell hardware via Amazon Braket. The stateful update
    prevents the memory collapse observed in stateless QRC beyond
    10 qubits.

    \index{stateful reservoir!implementation}

  \item \textbf{QHMM-OOM with forward-backward smoothing and TD($\lambda$)
    eligibility traces.} We formalized the QHMM as an Observable Operator Model,
    enabling exact trajectory probability computation and smoothing. The
    TD($\lambda$) trace mechanism allows credit assignment without
    backpropagation through the quantum processor, making online learning
    tractable.

    \index{TD($\lambda$)!contribution}
    \index{forward-backward!contribution}

  \item \textbf{Comprehensive benchmarks.} We evaluated QRC-EV on five
    datasets (synthetic and real-world EV data), demonstrating that the
    stateful QHMM achieves the best performance on memory-intensive tasks
    (NARMA-10, R²=0.86) while remaining competitive on periodic and
    chaotic tasks. On real EV charging data, classical Ridge regression
    (R²=0.91) significantly outperforms all quantum methods.

    \index{benchmark!comprehensive}
\end{enumerate}

\index{key result!QHMM NARMA-10}
\index{key result!Ridge EV dominance}

\section{Main Takeaways}

\index{takeaways}

The experimental results lead to three high-level conclusions:

\index{conclusions}

\begin{enumerate}
  \item \textbf{State is essential for quantum temporal learning.} The
    divergence between stateless QRC and stateful QHMM on NARMA-10
    (Figure~\ref{fig:narma10_scaling}) demonstrates that the depolarizing
    channel state maintenance is the critical ingredient enabling quantum
    reservoirs to solve tasks requiring long memory horizons. Without state,
    the effective memory capacity of a quantum circuit saturates and
    declines at relatively low qubit counts.

    \index{memory!stateful requirement}

  \item \textbf{Quantum advantage is real but narrow.} The QHMM-20q
    achieves R²=0.86 on NARMA-10, outperforming ESN (R²=0.52) and
    Ridge (R²=0.40) by 65\% and 115\% respectively. This is a meaningful
    advantage on a task that explicitly requires memory of 10 steps.
    However, this advantage does not transfer to the real-world EV charging
    task, where Ridge regression dominates. The conditions under which
    quantum temporal learning outperforms classical alternatives remain
    task-specific.

    \index{quantum advantage!narrow}
    \index{quantum advantage!conditions}

  \item \textbf{Method-task matching is critical.} For EV charging
    infrastructure operators, the choice between quantum and classical
    methods should be driven by the specific forecasting horizon and
    the underlying signal structure. For short-term load balancing,
    Ridge regression is sufficient and computationally efficient. For
    applications requiring detection of subtle temporal anomalies or
    multi-day demand forecasting, the QHMM's stateful modeling may
    provide marginal benefit at significant computational cost.

    \index{method selection!practical guide}
\end{enumerate}

\index{practical recommendations}

\section{Open Questions and Future Work}

\index{future work}

This work opens several directions for further investigation:

\index{future directions}

\begin{enumerate}
  \item \textbf{Physical quantum hardware evaluation.} All experiments
    reported here use CUDA Quantum simulation on NVIDIA GPUs. An important
    next step is to replicate these benchmarks on actual quantum processors
    accessible via Amazon Braket (e.g., IonQ Forte, Rigetti Aspen). The
    effect of gate errors and decoherence on QHMM performance remains
    an open empirical question.

    \index{physical quantum hardware}
    \index{Amazon Braket!physical QPUs}

  \item \textbf{Learned state discretization.} The current binning-based
    discretization of reservoir states into QHMM outcomes is a simple
    quantile approach that may discard useful state information. Future
    work should explore learned discretization methods, such as vector
    quantization or codebook learning, that adaptively partition the
    reservoir state space to maximize trajectory likelihood.

    \index{disambiguation!learned}
    \index{vector quantization!future}

  \item \textbf{Larger-scale EV forecasting.} The current evaluation
    is limited to the Palo Alto ACN network with 5000 hourly samples.
    Future work should evaluate on larger EV charging datasets spanning
    multiple networks, longer time horizons, and diverse geographical
    contexts. The relationship between network topology and QHMM structure
    is a particularly interesting direction.

    \index{generalization!larger datasets}

  \item \textbf{Adaptive TD($\lambda$) hyperparameters.} The trace decay
    parameters $\lambda_f$ and $\lambda_b$ were fixed at 0.8 throughout
    this work. An adaptive scheme that adjusts these parameters based on
    the observed TD error variance could improve credit assignment in
    non-stationary environments.

    \index{adaptive TD(λ)}

  \item \textbf{Integration with reinforcement learning for EV charging
    scheduling.} The QHMM's optimistic planning module produces action
    sequences that maximize expected cumulative reward. An end-to-end
    system that combines QHMM forecasting with RL-based charging
    infrastructure control (e.g., peak shaving, load balancing) would
    demonstrate practical utility.

    \index{reinforcement learning!charging scheduling}
    \index{peak shaving}
\end{enumerate}

\index{RL!EV charging}
\index{grid management}

\section{Closing Remarks}

\index{closing}

Quantum reservoir computing represents a promising direction for temporal
machine learning, particularly for tasks requiring explicit memory of
long histories. The QRC-EV framework demonstrates that the gap between
controlled quantum experiments and practical time-series forecasting is
narrowing—but not yet closed. The strongest practical advantage observed
in this work is on the NARMA-10 memory task, not on real-world EV
charging data. This underscores the importance of rigorous, task-specific
evaluation of quantum methods against well-tuned classical baselines.

\index{quantum ML!practical outlook}

The QHMM formalism provides a mathematically principled framework for
understanding the representational capacity of quantum dynamics for
temporal modeling. The forward-backward algorithm, eligibility traces,
and optimistic planning components are not specific to the CUDA Quantum
implementation—they apply to any quantum system whose dynamics can be
represented as CPTP channels. As quantum hardware improves, these methods
may become increasingly relevant.

\index{outlook!quantum ML}
\index{long-term outlook}

We hope this framework serves as a foundation for continued investigation
of quantum temporal learning in practical settings.

\index{reproducibility}
\index{open source}

\begin{itemize}
  \item \textbf{Code}: \url{https://github.com/rezacute/medseg3d-research}
  \item \textbf{Paper}: available as GitHub Actions artifact on branch
    \texttt{feature/agentic-research-workflow}
  \item \textbf{Results}: documented in \texttt{results/} with
    random seeds for reproducibility
\end{itemize}
